

\chapter{Questionario}
\newgeometry{a4paper, margin=1.5cm}
%\Large


\noindent{\bf Questionario \hspace{5cm} Codice:}
\medskip\hrule \medskip 

\noindent{\bf Età: \hspace{7cm} Classe:}
\newline
\newline
\noindent{\textbf{Genere:} $\bigcirc$ Femmina \quad $\bigcirc$ Maschio \quad $\bigcirc$ Altro \quad $\bigcirc$ Preferisco non rispondere }
\newline

\noindent\textbf{Istruzioni:} Immagina di trovarti nelle situazioni descritte qui sotto. Valuta ogni situazione in termini di quanta ansia proveresti durante l'evento specificato, mettendo una \textbf{crocetta} nella colonna che corrisponde al tuo grado di ansia. Ti ricordo che non ci sono risposte giuste o sbagliate.
\newline
\renewcommand{\arraystretch}{1.5}
\setlength{\tabcolsep}{4pt} % spazio orizzontale nelle celle


\noindent
\resizebox{\textwidth}{!}{%
\begin{tabular}{|p{7cm}|c|c|c|c|c|}
\hline
\multicolumn{1}{|c|}{\textbf{Situazione}} & \multicolumn{5}{c|}{\textbf{Grado di ansia}} \\
\cline{2-6}
\multicolumn{1}{|c|}{} & \textbf{Molto poca} & \textbf{Poca} & \textbf{Moderata} & \textbf{Abbastanza} & \textbf{Molta} \\
\hline
Usare le tabelline e gli schemi per svolgere esercizi di matematica &  &  &  &  &  \\
\hline
Pensare alla verifica scritta di matematica &  &  &  &  &  \\
\hline
Seguire con attenzione l'insegnante che risolve alla lavagna un'operazione di matemtaica difficile &  &  &  &  &  \\
\hline
Fare una verifica scritta di matematica &  &  &  &  &  \\
\hline
Svolgere per casa molti esercizi difficili di matematica &  &  &  &  &  \\
\hline
Seguire con attenzione la lezione di matematica &  &  &  &  &  \\
\hline
Seguire un compagno che risolve un esercizio di matematica &  &  &  &  &  \\
\hline
Essere interrogato “a sorpresa” in matematica &  &  &  &  &  \\
\hline
Affrontare un nuovo argomento di matematica &  &  &  &  &  \\
\hline
\end{tabular}}

\newpage

\noindent\textbf{Istruzioni:} Valuta quanto pensi di essere bravo/a ad affrontare le diverse situazioni descritte qui sotto. Metti una \textbf{crocetta} nella colonna che corrisponde al tuo grado di bravura. Ti ricordo che non ci sono risposte giste o sbagliate. 
\newline

\noindent\begin{tabular}{|p{7cm}|c|c|c|c|c|}
  \hline
  \textbf{Quanto sei bravo/a...} & \textbf{Per niente} & \textbf{Poco} & \textbf{Medio} & \textbf{Abbastanza} & \textbf{Molto} \\
  \hline
  Nell'imparare la matematica & & & & & \\
  \hline
  Nel risolvere problemi matematici & & & & & \\
  \hline
  Nel calcolo a mente & & & & & \\
  \hline
  Nel risolvere operazioni in colonna & & & & & \\
  \hline
  Nelle tabelline & & & & & \\
  \hline
  Nel finire i compiti che ti sono stati assegnati per casa & & & & & \\
  \hline 
  Nel concentrarti nello studio della matematica senza farti distrarre & & & & & \\
  \hline
  Impegnarti nello studio della matematica quando hai altre cose interessanti da fare & & & & & \\
  \hline
  Organizzarti nelle attività scolastiche & & & & & \\
  \hline
  Interessarti nelle materie scolastiche & & & & & \\
  \hline
\end{tabular}

\newpage

\noindent{\bf Questionario \hspace{5cm} Codice:}
\medskip\hrule \medskip 
\noindent\textbf{Istruzioni:} Valuta quanto sei d'accordo con le affermazioni proposte di seguito mettendo una \textbf{crocetta} nella casella che corrisponde al tuo pensiero. Ti ricordo che non ci sono risposte giuste o sbagliate. 
\newline

\noindent\begin{tabular}{|p{7cm}|c|c|c|c|c|}
  \hline
  \textbf{Quanto sei d'accordo?} & \textbf{Per niente} & \textbf{Poco} & \textbf{Medio} & \textbf{Abbastanza} & \textbf{Molto} \\
  \hline
  Mi piace studiare matematica & & & & & \\
  \hline
  La matematica è utile nella vita quotidiana & & & & & \\
  \hline
  Preferirei non avere lezioni di matematica & & & & & \\
  \hline
  Trovo la matematica interessante & & & & & \\
  \hline
  Penso che la matematica sia utile anche per il mio futuro & & & & & \\
  \hline
  La matematica è una materia che mi motiva a dare il meglio. & & & & & \\
  \hline 
  Non mi sento coinvolto/a durante le lezioni di matematica & & & & & \\
  \hline
  Vorrei approfondire la matematica oltre quello che facciamo a scuola & & & & & \\
  \hline
\end{tabular}


\newpage

\noindent\textbf{Istruzioni:} Valuta quanto sei d'accordo con le affermazioni proposte di seguito mettendo una \textbf{crocetta} nella casella che corrisponde al tuo pensiero. Ti ricordo che non ci sono risposte giuste o sbagliate. 
\newline

\noindent\begin{tabular}{|p{7cm}|c|c|c|c|c|}
  \hline
  \textbf{Quanto sei d'accordo?} & \textbf{Per niente} & \textbf{Poco} & \textbf{Medio} & \textbf{Abbastanza} & \textbf{Molto} \\
  \hline
  Sono preoccupato/a che i miei compagnini pensino che io abbia meno capacità matematiche a causa del mio genere & & & & & \\
  \hline
  Quando prendo un brutto voto in matematica temo che questo confermi lo stereotipo sul mio genere & & & & & \\
  \hline
  Penso che le ragazze non debbano occuparsi di matematica perchè non è adatta a loro & & & & & \\
  \hline
  Alcune volte penso che non dovrei essere bravo/a in matematica a causa del mio genere & & & & & \\
  \hline
  Mi preoccupo che la mia insegnante di matematica possa pensare che non sono bravo/a a causa del mio genere & & & & & \\
  \hline
  In matematica, i ragazzi hanno un talento naturale che le ragazze non possono compensare con l’impegno & & & & & \\
  \hline 
  Per ottenere gli stessi risultati in matematica i ragazzi devono impegnarsi di meno & & & & & \\
  \hline
  Le ragazze devono lavorare più dei ragazzi per ottenere gli stessi risultati in matematica & & & & & \\
  \hline
  Nonostante l’impegno, le ragazze restano mediamente meno brave in matematica dei ragazzi & & & & & \\
  \hline
\end{tabular}

































\restoregeometry
